\section*{Arbetslivserfarenhet}

\begin{tabular}{p{0.2\textwidth} | p{0.6\textwidth}}
    \cvevent{Jun 2025 - Nu}{Backend-utvecklare - Access Management}
    {\href{https://www.volvogroup.com/en/about-us/organization.html}{Volvo Group Digital Technology And Operations}}
    {Göteborg}{
        Utvecklar Identity-Providers och andra autentiseringskomponenter för telematikplattformen inom DTO:s Digital Services avdelning (tidigare känt som VGCS).
        Underhåller OAuth2-flöden med Java-mikrotjänster i Spring Boot på AWS.
    }{imgs/volvo.jpg} \\

    \cvevent{Sep 2023 - Jun 2025}{Mjukvaruingenjör - Innovation Lab}
    {\href{https://www.volvogroup.com/en/about-us/organization/volvo-group-connected-solutions.html}{Volvo Group Connected Solutions}}
    {Göteborg}{
        Byggde Lokate och andra innovationsprojekt med Spring Boot Java-mikrotjänster på AWS med Infrastructure as Code.
        Utvecklade TypeScript-lambdas samt frontends i React/React Native med GraphQL och WebSockets.
    }{imgs/volvo.jpg} \\

    \cvevent{Jun - Aug 2023}{IoT-utvecklare - Sommarjobbare}{
        \href{https://www.combitech.se/}{Combitech}}{Göteborg}{
        Skapade en Kotlin-middleware med coroutines för att koppla ihop ROS-baserade robotar med en IoT-plattform.
        Exponerade REST-API:er för åtgärder och implementerade MQTT/WebSockets för dataspårning i realtid.
    }{imgs/combitech.jpg} \\

    \cvevent{Jan - Jun 2023}{Examensarbete - VR \& Robot}{
        \href{https://www.combitech.se/}{Combitech}}{Göteborg}{
        Utvecklade en VR-applikation i Godot (C\#) för trådlös styrning av en ROS-baserad robot (Python/Ubuntu) via OpenXR.
    }{imgs/combitech.jpg}
\end{tabular}
